% Part: first-order-logic
% Chapter: sequent-calculus
% Section: rules-and-proofs

\documentclass[../../../include/open-logic-section]{subfiles}

\begin{document}

\iftag{FOL}
      {\olfileid{fol}{seq}{rul}}
      {\olfileid{pl}{seq}{rul}}

\olsection{Regels en \usetoken{P}{derivation}}

Hierna staan $\Gamma, \Delta, \Pi, \Lambda$ voor eindige rijen
!!{sentence}s.

\begin{defn}[Sequent]
Een \emph{sequent} is een uitdrukking van deze vorm:
\[
\Gamma \Sequent \Delta
\]
$\Gamma$ en $\Delta$ zijn eindige rijen !!{sentence}s van de taal
$\Lang L$. Elk van beide rijen mag leeg zijn. De rij links,
$\Gamma$, heet het \emph{antecedent}. De rij rechts, $\Delta$, heet het
\emph{succedent}.
\end{defn}

\begin{explain}
Een sequent geeft het volgende idee weer. Als alle !!{sentence}s in het
antecedent gelden, dan geldt ten minste een van de !!{sentence}s in het
succedent. Stel dat $\Gamma = \tuple{!A_1, \dots, !A_m}$ en $\Delta =
\tuple{!B_1, \dots, !B_n}$. Dan geldt $\Gamma \Sequent \Delta$ precies
als
\[
(!A_1 \land \cdots \land !A_m) \lif (!B_1 \lor \cdots \lor
!B_n)
\]
geldt. Er zijn twee bijzondere gevallen: $\Gamma$ is leeg of $\Delta$
is leeg. Als $\Gamma$ leeg is, dan is $m = 0$, en $\quad \Sequent
\Delta$ geldt precies als $!B_1 \lor \dots \lor !B_n$ geldt. Als
$\Delta$ leeg is, dan is $n = 0$, en $\Gamma \Sequent \quad$ geldt
precies als $\lnot(!A_1 \land \dots \land !A_m)$ geldt. Een sequent heet
geldig precies als de !!{sentence} die erbij hoort geldig is.
\end{explain}

Laat $\Gamma$ een rij !!{sentence}s zijn. Met $\Gamma, !A$ bedoelen we
dan de rij waarin $!A$ rechts achter~$\Gamma$ is gezet. Met $!A,
\Gamma$ bedoelen we de rij waarin $!A$ links voor~$\Gamma$ is gezet.
Als ook $\Delta$ een rij !!{sentence}s is, dan betekent $\Gamma,
\Delta$ dat we beide rijen achter elkaar zetten. De rij die zo ontstaat
heet hun concatenatie.

\begin{defn}[Beginsequent]
Een sequent heet een \emph{beginsequent} als hij
\iftag{prvFalse,prvTrue}{een van de volgende vormen heeft:
  \begin{enumerate}
    \item $!A \Sequent !A$
    \tagitem{prvTrue}{$\quad \Sequent \ltrue$}{}
    \tagitem{prvFalse}{$\lfalse \Sequent \quad$}{}
  \end{enumerate}}
{de vorm $!A \Sequent !A$ heeft}. Dit geldt voor elke !!{sentence} $!A$
van de taal.
\end{defn}

!!^{derivation}s in de sequentencalculus zijn bepaalde bomen van
sequenten. Bovenaan staan beginsequenten. Als een sequent onder een of
twee andere sequenten staat, moet hij daaruit correct volgen met een
afleidingsregel. De regels voor $\Log{LK}$ vallen in twee hoofdgroepen:
\emph{logische} regels en \emph{structurele} regels. Elke logische regel
is genoemd naar de !!{main operator} van de !!{sentence} met $!A$
en/of $!B$ in de onderste sequent. Zo'n regel heeft twee versies. De
ene versie leidt een sequent af waarin de !!{sentence} met de
!!{operator} links staat. Bij de andere versie staat die !!{sentence}
rechts.
\end{document}
